\section{Evaluation}
\label{sec:evaluation}

The evaluation answers three questions. First, does the refined VQE control
model satisfy the logical property suite? Second, do bug-seeded variants produce
counterexamples for representative protocol defects? Third, how does checking
cost change as the finite control bounds \maxiters{} and \ntermss{} increase?
All experiments were run with NuXmv 2.1.0 on the refined model
\texttt{vqe\_refined\_abort.smv}. The canonical configuration uses
\maxiters{} $=5$ and \ntermss{} $=3$.

\subsection{Canonical Model Checking Result}

Table~\ref{tab:canonical-results} summarizes the canonical verification run.
The refined model contains 41 checked properties: 10 CTL properties, 22 LTL
properties, and 9 invariant properties. NuXmv reports all 41 properties as
true. These properties cover reachability of successful and aborting terminal
states, control-flow conformance, data readiness, stale-data rejection,
terminal-state absorbency, and bounded termination.

\begin{table}[!t]
\centering
\caption{Canonical NuXmv verification result for the refined VQE control model.}
\label{tab:canonical-results}
\renewcommand{\arraystretch}{1.08}
\scriptsize
\begin{tabularx}{\columnwidth}{p{0.30\columnwidth}p{0.20\columnwidth}p{0.16\columnwidth}Y}
\toprule
\textbf{Property group} & \textbf{Logic} & \textbf{Count} & \textbf{NuXmv result} \\
\midrule
Reachability and termination & CTL & 10 & 10 true, 0 false \\
Control-flow and freshness & LTL & 22 & 22 true, 0 false \\
State consistency invariants & Invariant & 9 & 9 true, 0 false \\
\midrule
\textbf{Total} & CTL/LTL/INV & \textbf{41} & \textbf{41 true, 0 false} \\
\bottomrule
\end{tabularx}
\end{table}

The most important result is the separation between successful termination and
error termination. In the refined model, \done{} is reachable only through
successful convergence, whereas \abort{} is reachable for binding failure,
QPU or measurement failure, energy-computation failure, optimizer failure, and
non-convergence at the iteration bound. Thus the model does not silently treat
bounded non-convergence or backend failure as a successful VQE result.

\subsection{Counterexample-Based Mutation Tests}

To test whether the property suite detects meaningful logical faults, we checked
three intentionally faulty NuXmv variants. Each mutant changes one part of the
transition relation while keeping the same abstraction style. The expected
outcome is that NuXmv rejects the mutant and returns a counterexample trace.
Table~\ref{tab:mutation-results} summarizes the observed failures.

\begin{table*}[!t]
\centering
\caption{Mutation-test counterexamples produced by NuXmv.}
\label{tab:mutation-results}
\renewcommand{\arraystretch}{1.08}
\scriptsize
\begin{tabularx}{\textwidth}{p{0.07\textwidth}p{0.25\textwidth}p{0.31\textwidth}Y}
\toprule
\textbf{ID} & \textbf{Seeded logical fault} & \textbf{Violated representative property} & \textbf{Counterexample interpretation} \\
\midrule
M1 & Ignore a QPU or measurement error in \execmeasure{}. &
$\mathbf{G}((\texttt{state}=\execmeasure{} \wedge \texttt{qpu\_error}) \rightarrow \mathbf{X}(\texttt{state}=\abort{}))$ &
The trace reaches \execmeasure{} with \texttt{qpu\_error} true, but the model continues measuring, enters \computeenergy{}, and reaches \done{} instead of aborting. \\
\midrule
M2 & Keep stale \energyready{} true when returning to the next \bind{} state. &
$\mathbf{G}((\texttt{state}=\bind{} \wedge \texttt{iter}>0) \rightarrow \neg\energyready{})$ &
After \computeenergy{} and \optimizerstep{}, the next iteration reaches \bind{} with \texttt{iter}=1 while \energyready{} is still true, meaning the previous iteration's energy is incorrectly reusable. \\
\midrule
M3 & Enter \computeenergy{} before all Pauli terms have been measured. &
$\mathbf{G}(\texttt{state}=\computeenergy{} \rightarrow \termss{}=\ntermss{})$ &
With \ntermss{} $=3$, the trace reaches \computeenergy{} when \termss{} $=2$, demonstrating premature aggregation from an incomplete Hamiltonian-term measurement set. \\
\bottomrule
\end{tabularx}
\end{table*}

These counterexamples are useful because they show diagnostic value rather than
only positive verification. For M1, the trace exposes missing error handling at
the quantum-backend boundary. For M2, the trace exposes stale classical data
crossing an iteration boundary. For M3, the trace exposes premature aggregation
before all required Hamiltonian terms have been measured. Together, the three
mutants exercise the main classes of protocol errors that the model is intended
to rule out: missing abort behavior, stale-data reuse, and invalid phase
ordering.

Table~\ref{tab:ce-trace} gives a concrete excerpt from the M2 NuXmv trace. NuXmv
prints only variables whose values change at each step; the table expands the
unchanged values for readability. The violated formula is
$\mathbf{G}((\texttt{state}=\bind{} \wedge \texttt{iter}>0) \rightarrow
\neg\energyready{})$. The trace shows that after the first energy estimate is
computed, the mutant returns to \bind{} for iteration 1 while
\energyready{} is still true. This represents stale energy reuse across
iterations.

\begin{table}[!t]
\centering
\caption{Concrete NuXmv counterexample excerpt for stale energy reuse across iterations.}
\label{tab:ce-trace}
\renewcommand{\arraystretch}{1.08}
\scriptsize
\begin{tabularx}{\columnwidth}{p{0.13\columnwidth}p{0.27\columnwidth}p{0.13\columnwidth}p{0.18\columnwidth}Y}
\toprule
\textbf{Step} & \textbf{\texttt{state}} & \textbf{\texttt{iter}} & \textbf{\energyready{}} & \textbf{Interpretation} \\
\midrule
1.1 & \texttt{S\_INIT} & 0 & \false{} & Initial state. \\
1.2 & \bind{} & 0 & \false{} & First iteration starts with no energy estimate. \\
1.3 & \computeenergy{} & 0 & \false{} & Energy computation phase is reached. \\
1.4 & \checkstop{} & 0 & \true{} & Current iteration's energy is now valid. \\
1.5 & \optimizerstep{} & 0 & \true{} & Optimizer consumes the current energy. \\
1.6 & \bind{} & 1 & \true{} & Violation: next iteration starts while stale energy is still marked valid. \\
\bottomrule
\end{tabularx}
\end{table}

\subsection{Bounded Scaling Check}

Because this abstraction tracks the VQE control protocol rather than quantum
amplitudes, verification cost is primarily affected by the control bounds
\maxiters{} and \ntermss{}, not by Hilbert-space dimension. We therefore ran a
small bounded sweep over these two parameters. Table~\ref{tab:scaling-results}
reports local wall-clock time measured with \texttt{/usr/bin/time -p}. The times
are intended only as evidence that the finite control abstraction remains small
for the workshop-scale bounds.

\begin{table}[!t]
\centering
\caption{Bounded scaling results for the refined control abstraction.}
\label{tab:scaling-results}
\renewcommand{\arraystretch}{1.08}
\scriptsize
\begin{tabularx}{\columnwidth}{p{0.18\columnwidth}p{0.18\columnwidth}p{0.24\columnwidth}Y}
\toprule
\textbf{\maxiters{}} & \textbf{\ntermss{}} & \textbf{Properties true} & \textbf{Wall time} \\
\midrule
3 & 2 & 41/41 & 0.04 s \\
5 & 3 & 41/41 & 0.02 s \\
10 & 5 & 41/41 & 0.03 s \\
20 & 10 & 41/41 & 0.07 s \\
\bottomrule
\end{tabularx}
\end{table}

The sweep supports the intended scope of the model. It is not a quantum-state
scaling experiment and should not be interpreted as evidence about large-qubit
simulation. Instead, it shows that the protocol-level model checking remains
tractable as the bounded number of optimizer iterations and Hamiltonian-term
measurement groups increases within the finite abstraction.

\subsection{Evaluation Summary}

The evaluation supports three claims. First, the canonical refined VQE control
model satisfies the complete 41-property suite. Second, the property suite
rejects representative bug-seeded variants and returns short, interpretable
counterexample traces. Third, the bounded control abstraction is inexpensive to
check for the parameter settings considered here. These results establish the
intended workshop contribution: a machine-checkable logical baseline for the VQE
execution protocol, including explicit distinction between successful
termination and safety-preserving abort behavior.
